\originalchapter{树与生成树}\label{ch:trees}
\sourceof{6.042J 第11.10节; 18.212 L22图论部分、L23}
\originalsection{树的等价刻画}
树是连通而没有圈的图. 它没有多余的连接, 又能连接全部顶点, 所以在网络建设、递归结构和计数中都自然出现. 一个没有圈的图称为森林, 每个分支都是树; 单顶点也算树.
\begin{lemma}\label{lem:leaf}
含至少两个顶点的有限树至少有两个度1顶点, 称为叶子.
\end{lemma}
\begin{proof}
取最长路 $v_0\cdots v_\ell$. 端点 $v_0$ 若有路外邻居就能延长路; 若有除 $v_1$ 之外的路内邻居就产生圈. 因而它只有一个邻居. 另一端点同理.
\end{proof}
\begin{theorem}\label{thm:tree-equivalent}
对 $n\ge1$ 的简单图 $G$, 以下条件等价: (1) $G$ 是树; (2) 任意两顶点之间恰有一条路; (3) $G$ 连通且 $m=n-1$; (4) $G$ 无圈且 $m=n-1$; (5) $G$ 连通, 删除任意边都不连通.
\end{theorem}
\begin{proof}
树中两顶点有路. 两条不同的路沿首次分开再会合的区段产生圈, 所以路唯一. 唯一路的任一边删除后, 其两端不能再连通, 得到(5). 反向, 若连通图有圈, 删除圈上一边仍可沿圈另一侧通行, 违反(5).

由叶子引理, 对树删去一片叶子得到小一阶的树, 从单顶点归纳得到 $m=n-1$. 任意连通图都可不断删除圈上的边而保持连通, 直到得到一棵包含全部顶点的树, 称为生成树. 它有 $n-1$ 条边. 若原图本来只有 $n-1$ 条边, 就无可删除的边, 原图本身为树, 得到(3)的反向. 森林若有 $c$ 个分支, 在各分支分别计数得到 $m=n-c$; (4)使 $c=1$. 这同时完成全部等价关系.
\end{proof}
树加一条新边, 恰好产生一个圈: 它由新边与树中连接其端点的唯一路组成. 树删一条边恰好分成两个分支: 端点已不连通, 而每个顶点沿原来的路总能到达至少一个端点. 后面的最小生成树算法将围绕这两个操作构造交换论证.
\begin{example}
一棵树中度1、2、3的顶点个数分别为 $a,b,c$, 且没有其他度数. 求 $a$ 与 $c$ 的关系.
\end{example}
\begin{solution}
若树有至少两个顶点, 握手引理给出 $a+2b+3c=2(a+b+c-1)$, 所以 $a=c+2$. 度2顶点只把边细分, 不影响分枝与叶子之间的数量关系. 单顶点度0, 已不满足题面所列度数种类.
\end{solution}
\originalsection{根树与生成森林}
指定根后, 每个非根顶点在通向根的唯一路上有一个父顶点, 其余相邻顶点为子顶点. 顶点的深度是到根的距离. 从根向外逐层访问得到的父边必组成生成树, 但不同访问策略可能得到不同树.

在一般图中, 极大的无圈边集必在每个原分支内连通: 若一个原分支被分为两块, 可加入跨两块的边而不产生圈, 与极大性矛盾. 因而它是生成森林, 有 $n-c(G)$ 条边. 注意“极大”是不能再包含更多边, “最大”是边数达到最大; 这里因为森林的结构, 二者最终边数相同, 不能在其他优化问题中随便混用.
\begin{center}
\begin{tikzpicture}[x=1.25cm,y=1cm]
\node[vertex](r)at(0,0){$r$};\node[vertex](a)at(-1,-1){$a$};\node[vertex](b)at(1,-1){$b$};
\node[vertex](c)at(-1.6,-2){$c$};\node[vertex](d)at(-.4,-2){$d$};\node[vertex](e)at(1,-2){$e$};
\foreach \u/\v in {r/a,r/b,a/c,a/d,b/e}\draw[edge](\u)--(\v);
\draw[dashed,bend right=35](c)to(e);
\node at(3,-.6){实线构成生成树};\node at(3,-1.2){加虚线产生唯一圈};
\end{tikzpicture}
\end{center}
图中虚线若加入, 圈为 $c,a,r,b,e,c$. 从这条圈中删去任一原树边, 就得到另一棵生成树. 没在圈上的边不能代替它被删, 因为那样仍有圈且另有部分被断开.
\originalsection{Prüfer 编码与 Cayley 公式}
\begin{theorem}[Cayley 公式]\label{thm:cayley}
顶点集为 $[n]=\{1,\ldots,n\}$ 的树共有 $n^{n-2}$ 棵, 其中 $n\ge2$.
\end{theorem}
\begin{proof}
给每棵标号树作 Prüfer 编码: 每次删除当前编号最小的叶子, 记录其邻居编号, 直到只剩两个顶点. 得到长度 $n-2$ 的序列 $(a_1,\ldots,a_{n-2})$, 每项在 $[n]$ 中.

反向给定任意这种序列. 在尚未删除的标号中, 取未在当前剩余序列出现的最小标号 $b_i$, 加边 $b_i a_i$, 删去 $b_i$ 并删掉首项 $a_i$. 该标号存在, 因为当前顶点数比序列长度多2. 序列中出现的标号不会提前被删除, 所以 $a_i$ 是一个仍在的、不同于 $b_i$ 的顶点. 最后连接剩下的两顶点.

构造出的图连通: 每个删除顶点都连到一个以后才删除的顶点, 顺着删除次序最终到达最后一条边. 它有 $n-1$ 条边, 因而是树. 在当前构造树中, 一个顶点的度数是它在当前序列中的出现次数加1: 每次作为 $a_i$ 接收一条边, 最后作为 $b_i$ 或最后两点之一又有一条边. 所以当前叶子恰是序列中不出现的标号, 两个过程每一步选出的最小叶子一致. 编码与解码互逆, 树数等于序列数 $n^{n-2}$.
\end{proof}
\begin{corollary}
若正整数 $d_1,\ldots,d_n$ 满足 $\sum d_i=2n-2$, 则具有该指定度数的标号树共有
\[
\frac{(n-2)!}{\prod_{i=1}^n(d_i-1)!}.
\]
\end{corollary}
\begin{proof}
由编码证明, 标号 $i$ 在 Prüfer 序列中恰出现 $d_i-1$ 次. 这些次数之和为 $n-2$, 按多重排列计数即得. 条件也使这样的序列确实存在.
\end{proof}
\begin{example}
解码序列 $(4,4,2,4)$, 并求树的度数.
\end{example}
\begin{solution}
这里 $n=6$. 首先缺失的最小标号为1, 加边 $14$; 剩余序列 $(4,2,4)$ 中缺失的最小尚存标号为3, 加边 $34$. 然后对 $(2,4)$ 取5, 加边 $52$; 对 $(4)$ 取2, 加边 $24$; 最后连接4、6. 度数为 $(1,2,1,4,1,1)$, 与每个标号出现次数加1一致.
\end{solution}
\originalsection{习题与解答}
\begin{exercise}\label{ex:tree1}
证明任意树都是二部图, 并说明它何时存在完美匹配.
\end{exercise}
\begin{exercise}\label{ex:tree2}
顶点为 $[n]$, 指定顶点1的度恰为 $k$, 其中 $1\le k\le n-1$. 求树的数量.
\end{exercise}
习题\ref{ex:tree1}的解答.
\begin{solution}
树没有圈, 特别没有奇圈, 由定理\ref{thm:bipartite}为二部图. 完美匹配可以从叶子判断: 叶子 $u$ 的唯一边 $uv$ 必须被使用, 删去 $u,v$ 后, 原树有完美匹配当且仅当剩余森林有完美匹配. 反复执行, 最终为空图则存在, 出现孤立顶点则不存在. 偶数阶仅是必要条件, 例如四叶星有五点, 而五叶星有六点也无完美匹配. 匹配的正式理论见第\ref{ch:matching}章.
\end{solution}
习题\ref{ex:tree2}的解答.
\begin{solution}
在长度 $n-2$ 的序列中选 $k-1$ 个位置填1, 其余位置任填另外 $n-1$ 个标号. 数量为 $\binom{n-2}{k-1}(n-1)^{n-k-1}$, 边界 $k=n-1$ 时只有全1序列, 对应以1为中心的星.
\end{solution}
